\documentclass[../../../include/open-logic-section]{subfiles}

\begin{document}

\olfileid[pt]{sfr}{infinite}{card-sb}

\olsection{Apêndice: Provando Schröder-Bernstein}

Antes de deixarmos a teoria ingênua dos conjuntos, há ainda uma última
prova ingênua (mas sofisticada!) a considerar. Trata-se de uma prova de
Schr\"oder-Bernstein (\olref[sfr][siz][sb]{thm:schroder-bernstein}): se
$\cardle{A}{B}$ e $\cardle{B}{A}$ então $\cardeq{A}{B}$; isto é, dados
!!{injection}s $f \colon A \to B$ e $g \colon B \to A$, existe
!!a{bijection} $h \colon A \to B$.

Neste capítulo, seguimos a noção dedekindiana de \emph{fechos}. De
fato, Dedekind deu uma prova muito elegante de Schr\"oder-Bernstein com
esta noção, e aqui a apresentamos. A prova segue de perto
\citet[pp.~157--8]{Potter2004}, se quiser um tratamento ligeiramente
diferente mas essencialmente semelhante. Uma breve pesquisa na Internet
também convencerá o leitor de que este é um teorema---tal como a
irracionalidade de $\sqrt{2}$---para o qual existem \emph{muitas}
provas interessantes e distintas.

Usando notação semelhante à de \olref[sfr][infinite][dedekind]{Closure},
seja
\[
\Closureofunder{f}{B} = \bigcap \Setabs{X}{B \subseteq X 
\text{ e $X$ é $f$-fechado}}
\]
para cada conjunto $B$ e função~$f$. Assim definido,
$\Closureofunder{f}{B}$ é o menor conjunto $f$-fechado que contém~$B$,
no sentido de que:

\begin{lem}\ollabel{Closureprops}
	Para qualquer função $f$, e qualquer $B$:
	\begin{enumerate}
		\item\ollabel{Closurehaselem} $B \subseteq \Closureofunder{f}{B}$; e
		\item\ollabel{Closureclosed} $\Closureofunder{f}{B}$ é $f$-fechado; e
		\item\ollabel{Closuresmallest} se $X$ é $f$-fechado e $B
		\subseteq X$, então $\Closureofunder{f}{B} \subseteq X$.
	\end{enumerate}
\end{lem}

\begin{proof}
Exatamente como no \olref[sfr][infinite][dedekind]{closureproperties}.
\end{proof}

Precisamos de mais um fato para chegar a Schr\"oder-Bernstein:

\begin{prop}\ollabel{sbhelper}
Se $A \subseteq B \subseteq C$ e $A \approx C$, então $\cardeq{\cardeq{A}{B}}{C}$.
\end{prop}

\begin{proof}
Dada !!a{bijection} $f \colon C \to A$, seja $F =
\Closureofunder{f}{C \setminus B}$ e defina-se uma função $g$ com
domínio $C$ do seguinte modo:
\[
	g(x) = 
	\begin{cases}
			f(x) &\text{se $x \in F$}\\
			x & \text{caso contrário}
		\end{cases}
\]
Mostraremos que $g$ é !!a{bijection} de $C$ para $B$, donde seguirá que
$\comp{f^{-1}}{g} \colon A \to B$ é !!a{bijection}, completando a
prova.

Primeiro alegamos que, se $x \in F$ mas $y\notin F$, então $g(x) \neq
g(y)$. Por reductio, suponhamos o contrário, de modo que $y = g(y) = g(x) =
f(x)$. Como $x \in F$ e $F$ é $f$-fechado por \olref{Closureprops},
temos $y = f(x) \in F$, contradição.

Suponhamos agora $g(x) = g(y)$. Então, pelo que precede, $x \in F$ sse
$y \in F$. Se $x, y \in F$, então $f(x) = g (x) = g(y) = f(y)$, logo
$x = y$ pois $f$ é !!a{bijection}. Se $x, y \notin F$, então $x = g(x) =
g(y) = y$. Logo $g$ é !!a{injection}.

Resta mostrar que $\ran{g} = B$. Fixemos $x \in B \subseteq C$. Se $x
\notin F$, então $g(x) = x$. Se $x \in F$, então $x = f(y)$ para algum
$y \in F$, pois caso contrário $F \setminus \{x\}$ seria $f$-fechado e
estenderia $C\setminus B$, o que é impossível por \olref{Closureprops};
agora $g(y) = f(y) = x$.
\end{proof}

Finalmente, eis a prova do resultado principal. Recordemos que, dada uma
função $h$ e um conjunto $D$, definimos $\funimage{h}{D} = \Setabs{h(x)}{x
\in D}$.

\begin{proof}[Prova de Schr\"oder-Bernstein] Sejam $f \colon A \to B$
e $g \colon B \to A$ !!{injection}s. Como $\funimage{f}{A}
\subseteq B$, temos $\funimage{g}{\funimage{f}{A}} \subseteq
g[B] \subseteq A$. Além disso, $\comp{f}{g} \colon A \to
\funimage{g}{\funimage{f}{A}}$ é !!a{injection}, pois tanto $g$ como
$f$ o são; e de fato $\comp{f}{g}$ é !!a{bijection}, apenas pela forma
como definimos o seu contradomínio. Logo
$\cardeq{\funimage{g}{\funimage{f}{A}}}{A}$, e portanto, por
\olref{sbhelper}, existe !!a{bijection} $h \colon A \to
\funimage{g}{B}$. Além disso, $g^{-1}$ é !!a{bijection}
$\funimage{g}{B} \to B$. Logo $\comp{h}{g^{-1}} \colon A \to B$ é
!!a{bijection}.
\end{proof}

\end{document}
